संच~$A$ वरील \emph{संबंध}~$R$ हा $A$ च्या घटकांना एकमेकांशी जोडण्याचा एक मार्ग आहे. $x$ आणि~$y$ यांच्यात तो संबंध \emph{लागू असेल}, तर आपण $Rxy$ असे लिहितो. आकारिक दृष्टीने, ज्या $\tuple{x,y} \in A^2$ जोड्यांसाठी $Rxy$ लागू आहे त्यांचा संच म्हणजे $R$ असे मानता येते. लहान असणे, मोठे असणे, समान असणे, निःशेष भाग देणे, समान लांबीचे असणे, उपसंच असणे आणि समान आकारमानाचे असणे ही (संख्यांच्या, चिन्हमालांच्या किंवा संचांच्या संचांवरील) संबंधांची महत्त्वाची उदाहरणे आहेत. संबंध दृश्य स्वरूपात दाखवण्याचा एक सर्वसाधारण मार्ग म्हणजे \emph{आलेख}. आलेखाकडे एक द्विपदी संबंध (\emph{कड} संबंध) आणि त्याच्या आधारभूत \emph{शिखरांचा} संच यांचा मिळून बनलेला घटक म्हणूनही पाहता येते.

काही संबंधांतील समान गुणधर्म त्यांना विशेष रंजक किंवा उपयुक्त बनवतात. प्रत्येक घटकाचा स्वतःशी $R$ संबंध असेल, तर संबंध~$R$ \emph{परावर्ती} असतो. कोणत्याही $x$ आणि $y$ साठी $Rxy$ लागू असताना $Ryx$ देखील लागू असेल, तर तो \emph{सममित} असतो. $Rxy$ आणि $Ryz$ या दोन्हींवरून $Rxz$ निश्चित होत असेल, तर तो \emph{संक्रमक} असतो. हे तिन्ही गुणधर्म असलेल्या संबंधांना \emph{तुल्यता संबंध} म्हणतात. $Rxy$ आणि $Ryx$ या दोन्हींवरून $x=y$ निश्चित होत असेल, तर संबंध \emph{प्रतिसममित} असतो. जे संबंध परावर्ती, प्रतिसममित आणि संक्रमक असतात ते \emph{अंशतः क्रम} असतात. कोणत्याही $x$ आणि~$y$ साठी $Rxy$, $x=y$ किंवा $Ryx$ यांपैकी किमान एक लागू होतो असा अंशतः क्रम म्हणजे \emph{रेषीय क्रम}. (सर्वसाधारणपणे या गुणधर्माचा संबंध \emph{संयोजित} असतो.)

संबंध हे (जोड्यांचे) संच असल्याने त्यांच्यावर संचांच्या क्रिया करता येतात (उदाहरणार्थ, संबंधांचा संयोग आणि छेद घेता येतो). संबंध एकापुढे एकही जोडता येतात (\emph{सापेक्ष गुणाकार}~$R \mid S$). $R$ च्या सर्व धन घातांचा संयोग म्हणजे त्याचे \emph{संक्रमक संवरण}~$R^+ = \bigcup_{0 < n \in \Nat} R^n$. येथे $R^1 = R$ आणि $R^{n+1} = R^n \mid R$.
